\section{Supplementary Results on the Diagonal-Equation Solvers}
\label{app:sde-side}

\noindent\textbf{Status.} AI-checked, pending human review; the barrier remark is argued, not proved; provenance in Appendix~\ref{app:provenance}, items~\ref{prov:sde-is-runtime}, \ref{prov:sde-threshold} and \ref{prov:sde-subquadratic}--\ref{prov:sde-worst-case}.

This appendix collects material of Section~\ref{sec:schedules} that is not needed for the statements of the schedules: the operation counts and the tape layout behind the Turing-machine bounds of Remark~\ref{rem:is-runtime}, the barriers that stand between the proved two-level threshold and the square gap of Remark~\ref{rem:two-level-threshold-size}, and three remarks on the arity of the diagonal equation below the two-level count.

\subsection{Tape layout and run counts for the Schedule 1 solver}
\label{app:sde-side-tape}

\begin{remark}[Operation counts, tape layout and run counts of Schedule~1; AI-checked, pending human review; Turing-machine bounds sketched, constants not computed]
\label{rem:sde-side-tape}
Keep the notation of Remark~\ref{rem:is-runtime}: $\kappa = \kappa_d$, $a = g_d$, $b = \lceil\log_2 p\rceil$, $\mathcal{G} = G(\kappa,a)$, $\nu$, $\pi$, $H_\pi$ and the $\pi$- and $\pi'$-parts.

\emph{The variant.} The following deterministic variant is defined for every $\kappa \ge 2$; when $\kappa = q^e$ is a prime power it makes at most $e \le \log_2\kappa$ runs, and a single run, never repeated, when $\kappa$ is prime, and for general $\kappa$ it makes at most $r_{\mathrm{run}}$ runs (AI-checked, pending human review). Before the first run it computes the kernel vector $\alpha$ of the first $a+1$ columns, which is the kernel of the first node of \cite[Lemma~12]{IvanyosSantha2015} and does not depend on the classifier, and it starts the classifier at a generator $\eta_1$ of the group $H_1$ generated by the $\pi$-parts of the nonzero $\alpha_k$. Whenever a restart enlarges $\eta$, it first checks whether these $\pi$-parts lie in one coset of $\langle\eta\rangle^{\kappa}$. If they do, then, since the $\pi'$-parts are $\kappa$-th powers, discrete logarithms in $\langle\eta\rangle$ write $\alpha_k = \gamma_0 u_k^{\kappa}$, with $u_k \ne 0$ exactly when $\alpha_k \ne 0$, so $\sum_k u_k^{\kappa}\lambda_k = 0$ is a nonzero witness on the first $a+1$ columns, and the variant stops. If $\kappa = q^e$, each restart multiplies $|\langle\eta\rangle|$ by at least $q$, so after $e$ of them $|\langle\eta\rangle| \ge \kappa|H_1|$; subgroups of the cyclic $q$-group are totally ordered, hence $H_1 \subseteq \langle\eta\rangle^{\kappa}$ and the check succeeds. A generator of all of $H_\pi$ cannot be fixed in advance instead: for suitable columns every completed run uses one, and it is a $q$-th-power nonresidue for every prime $q \mid \kappa$, which no known deterministic polynomial-time method constructs unconditionally (under the generalized Riemann hypothesis one is found by search).

\emph{Operation counts.} Each run makes $\sum_{i<L}\mathcal{G}^i$ calls of \cite[Theorem~11]{IvanyosSantha2015}, with $L = \lceil\log_2 \ell_d\rceil$. One call runs \cite[Lemma~12]{IvanyosSantha2015} with $\mathcal{G}/\mathcal{G}_j$ nodes at level $j$, where $\mathcal{G}_j = \kappa^{j(j-1)/2}(a+1)^j$. Each node is one kernel computation for a $\kappa^{j-1}a\times\kappa^{j-1}(a+1)$ matrix and at most $\kappa^{j-1}(a+1)$ classifications, and each classification is a Pohlig--Hellman logarithm of $O((\nu+\log\kappa)b+\kappa\log^2\kappa)$ multiplications. The zero sum of $\kappa$-th powers, computed by the reduction of \cite[Section~6.1]{vandeWoestijne2005} in the variant used by the implementation of Section~\ref{sec:schedule1} (Appendix~\ref{app:provenance}, item~\ref{prov:sde-is-solver}), uses $O(\kappa^3 b^2\log\kappa + \kappa^2\nu\, b^2\log\kappa)$ field operations, hence $O(\kappa^3 b^2\log\kappa + \kappa^2 b^3\log\kappa)$ since $\nu \le b-1$, independently of $a$ and of the columns. Counting all of this, together with the Bezout step of Remark~\ref{rem:is-import}(b), gives at most $40\kappa^3 N_{\mathrm{col}} b(a^2+b^2)$ multiplications and inversions in $\Fp$ and at most twice as many additions, where $N_{\mathrm{col}} = \mathcal{G}^{L} \le B_{\kappa}(a) \le M_d$; at $\kappa_d = 1$ the count is $a^2(a+1)+2a+2b(a+1)$. With schoolbook arithmetic and the binary extended Euclidean algorithm, each of these operations costs $O(b^2)$ steps of a multitape Turing machine that keeps its operands and $p$ on dedicated tapes.
% NEEDS HUMAN VERIFICATION (referee 9, 2026-11-04): the ISSAC 2005 paper has six pages (pp. 348--353); check whether "Section~6" / "Section~6.1" of vandeWoestijne2005 refers to it or to the thesis of the same title.

\emph{Tape layout.} All other work is sequential data movement under the following layout (AI-checked, pending human review). The tree of \cite[Lemma~12]{IvanyosSantha2015} is evaluated level by level, for all calls of one accumulation level of \cite[Section~3.3]{IvanyosSantha2015} together. The children of a level-$j$ node are $\kappa^{j-1}(a+1)$ consecutive records of level $j-1$, so their concatenation is the node's matrix stored by columns, and each Gauss--Jordan pivot is one sweep over it, with $O(b)$ head movement per entry against one multiplication per entry. The vectors $w_{a'}$ are accumulated in one sweep per child over $\kappa^j$ slots. The inputs below a node, and the columns below a call, occupy consecutive positions, so the classes, the coefficients and the representation indices are updated by merging two streams written in the same order. The Pohlig--Hellman digit tables, the collision table of the zero sum and the values $z_j$ have at most $\kappa+1$ entries and are scanned linearly, and a restart erases the work tapes at no more than the cost of the aborted run. Counted sweep by sweep, the head movement exceeds a constant per field operation only by factors polynomial in $\kappa$, and it stays below an absolute constant times $\kappa^3N_{\mathrm{col}}b^2(a^2+b^2)$ per run. The order in which the nodes are evaluated affects neither correctness nor the bounds $\nu+1$, $e$ and $r_{\mathrm{run}}$ (below) on the number of runs; the kernel vector $\alpha$ of the variant is computed before the first run in either order. With the $O(b^2)$ cost of one field operation this gives the bounds $T^{\mathrm{IS}}_d$ and $T^{\mathrm{IS},\ast}_d$ of Remark~\ref{rem:is-runtime}.

\emph{Run counts.} For every $\kappa_d \ge 2$ the variant makes at most $r_{\mathrm{run}} = \nu - \min_{q \mid \kappa_d} v_q\bigl((p-1)/\kappa_d\bigr)$ runs, with $q$ prime and $v_q$ the $q$-adic valuation; this is $e$ when $\kappa_d = q^e$, and $1 + \max_{q \mid \kappa_d} v_q(p-1) \le r_{\mathrm{run}} \le \nu$ when $\kappa_d$ has at least two distinct prime factors (proved; for anchors with at most three coefficients, an exhaustive search over subgroup chains at $\kappa_d \in \{6,10,12,15,30\}$ and $p < 4000$ gives the same maximum, a finite computation). When $(\kappa_d(g_d+1)-1)\,g_d \ge \nu-1$, explicit columns make the variant use exactly $r_{\mathrm{run}}$ runs and the algorithm of Theorem~\ref{thm:is-import} exactly $\nu+1$ runs; when $\kappa_d$ has at least two distinct prime factors, other explicit columns make the variant use exactly $r_{\mathrm{run}}$ runs for every $g_d \ge 1$ provided $\nu \le 10^{12}$ (proved for the default kernel vector of the implementation, under depth-first and under level-by-level evaluation; no full run at composite $\kappa_d$ was executed). Hence, when $\kappa_d$ has at least two distinct prime factors, no bound on the number of runs in terms of $\kappa_d$ alone holds for either algorithm, the run count cannot remove the factor $b$, and the bound $T^{\mathrm{IS}}_d$ stands; these are counts of runs, not lower bounds on time (AI-checked, pending human review). The derivation of the operation counts, an instrumented check of them on the implementation at $\kappa_d\le 3$, and the tests of the variant are recorded in Appendix~\ref{app:provenance}, item~\ref{prov:sde-is-runtime}.
\end{remark}

\subsection{Barriers for the two-level threshold}
\label{app:sde-side-barrier}

\begin{remark}[Barriers for the two-level threshold; argued; AI-checked, pending human review]
\label{rem:sde-side-barrier}
Keep the notation of Remark~\ref{rem:two-level-threshold-size}. The following arguments are made in the sources of Appendix~\ref{app:provenance}, item~\ref{prov:sde-threshold}; they are argued, not proved here. By them, closing the square gap needs information that the standard tools do not supply. If $p \le \kappa^\kappa/2$, no argument that averages over $\beta$ and uses only statements about at most $\kappa-1$ of the column counts $|X_j(\beta)|$ at a time, each true when the columns are independent, can prove $p \notin E_2(\kappa)$: there is a law of the column counts that agrees with independence on every $\kappa-1$ columns and gives probability $0$ to all $\kappa$ columns being nonempty (Benjamini, Gurel-Gurevich and Peled~\cite{BenjaminiGurelGurevichPeled2012}, Proposition~26), and exact searches find primes, such as $p = 37, 53$ for $\kappa = 4$ and $p = 131$ for $\kappa = 5$, where any $\kappa-1$ columns can be filled simultaneously but no contraction with every entry nonzero exists. Bounding the full character expansion of the number of contractions term by term cannot certify any $p \le 2^{-1/(\kappa-1)}\kappa^\kappa$, because the bound applies to every coset pattern at once and the configurations with all $\alpha_j$ equal overfill the row-constant patterns; and the linear forms $\alpha_j + \beta_i$ are not in general position, so nonsingular Deligne--Katz estimates do not apply. Treating the degenerate configurations exactly does not change this picture. Every contraction counted lies in the stratum where the $\alpha_j$ are distinct and the $\beta_i$ are distinct, and on that stratum no coset pattern is overfull; but inclusion--exclusion moves the degenerate mass into $\kappa^2(\kappa-1)^2$ one-line character sums and the $2\times2$-block sums, each of size about $p^{2\kappa-1}$, while the roughly $\kappa^{\kappa^2-1}$ sums that involve every row and column are at best of square-root size $p^{\kappa+1}$, a bound no known estimate supplies. With the exact value of every sum, the refined term-by-term certificate first holds at $p = 97 > T_A(3)$ for $\kappa = 3$; keeping the low-order sums with their phases lowers this to $p = 43$, but at $\kappa = 4$ the same hybrid certificate still fails at $p = 1009 > T_A(4)$ in a sampled evaluation, and for larger $\kappa$ such certificates can improve on $T_A(\kappa)$ by at most a factor tending to $1$, far from $(\kappa-1)^{\kappa+1}$ (finite computations, AI-checked, pending human review). In the other direction, a vanishing cyclotomic number gives a two-row obstruction, so every prime $p \equiv 1 \pmod\kappa$ with $p \le \kappa^2 - 3\kappa + 3$ lies in $E_2(\kappa)$, and $E_2(D) \subseteq E_2(\kappa)$ on primes $p \equiv 1 \pmod\kappa$ for every divisor $D \ge 2$ of $\kappa$ (Appendix~\ref{app:s2-local}; proved, AI-checked, pending human review); obstructions confined to $r$ rows are void above $((r-1)\kappa^r+2)^2$, so for fixed $r$ they cannot certify primes near $\kappa^{(\kappa+1)/2}$ once $\kappa$ is large, and the square gap stays open. A sub-family of the problem is a coloured-clique question in $\Fp$ of the same type as the clique number of the Paley graph, where the analogous gap is a long-standing open problem.
\end{remark}

\subsection{Side results on the arity of the diagonal equation}
\label{app:sde-side-arity}

The three remarks of this subsection concern the arity of the diagonal equation~\eqref{eq:sde-target} below the two-level count; they are referred to from Remark~\ref{rem:triangular-scope}, and Remark~\ref{rem:sde-block-steered} supplies the block-steered option of Remark~\ref{rem:combined-schedule}.

\begin{remark}[Below quadratic arity; AI-checked, pending human review]
\label{rem:sde-subquadratic}
This remark is AI-checked, pending human review (Appendix~\ref{app:provenance}, item~\ref{prov:sde-subquadratic}); no human has checked it, and nothing else in the paper uses it. Keep the notation of Proposition~\ref{prop:triangular-sde}. For a set $S$ of columns and a subspace $U\le\Fp^{a}$ put $R_U(S)=\{\sum_{i\in S}x_i\lambda_i+U : x\in(P\cup\{0\})^{S}\}\subseteq\Fp^a/U$. Call $U$ \emph{realizable} by disjoint sets of columns if every element of $U$ is a value $\sum_i x_i\lambda_i$ with weights $x_i\in P\cup\{0\}$ on their union. If $U$ is realizable and $S$ is disjoint from the sets used, then every subspace $V\supseteq U$ with $V/U\subseteq R_U(S)$ is realizable, as the spaces $U_{k+1}=U_k\oplus\Fp v_{k+1}$ are in the proof of Proposition~\ref{prop:triangular-sde}.

\textup{(a)} \emph{One dependency gives one direction.} Suppose the images of $S$ in $\Fp^a/U$ form a circuit: they are dependent and every proper subset is independent, so their relations are the multiples of one $\delta\in\Fp^S$ with every $\delta_i\ne0$. Suppose also that the entries of $\delta$ do not all lie in one coset $g^\nu P$. Then $R_U(S)$ contains no $2$-dimensional subspace. This is a property of the set $R_U(S)$, so it bounds every rule for choosing the weights, universal or depending on the columns, whatever its running time. Indeed, choose $i,j\in S$ with $\mathrm{cos}(\delta_i)\ne\mathrm{cos}(\delta_j)$ and let $W_0$ be the span of the images of $S\setminus\{i,j\}$, a hyperplane in the span $W$ of the images of $S$. A plane $T\subseteq R_U(S)$ lies in $W$, so it contains some $h\ne0$ in $W_0$. Let $z\in\Fp^S$ be the vector supported on $S\setminus\{i,j\}$ with $\sum_l z_l\lambda_l+U=h$; it is unique and nonzero. For $\gamma\in\Fp^\times$ the vectors $x\in\Fp^S$ with $\sum_l x_l\lambda_l+U=\gamma h$ are $\gamma z+s\delta$, $s\in\Fp$. For $s\ne0$ their $i$th and $j$th entries $s\delta_i$ and $s\delta_j$ lie in different cosets, so such a vector lies in $(P\cup\{0\})^S$ only if $s=0$ and $\gamma z\in(P\cup\{0\})^S$, which fixes $\mathrm{cos}(\gamma)$. Since $\kappa\ge2$, some multiple $\gamma h$ is not in $R_U(S)$, a contradiction. The hypothesis on $\delta$ is needed: for $\kappa=2$ and $p\in\{29,53,101\}$ the columns $e_1,e_2,-e_1-e_2$, with $\delta=(1,1,1)$, have $R_0(S)=\Fp^2$; such a $\delta$, however, is itself a solution. Consider a scheme that enlarges the realizable space one set at a time, each set a circuit modulo the space realized so far, and finishes when some set realizes an element of that space with weights not all zero. A circuit whose dependency meets two cosets adds at most one dimension, and it cannot finish the scheme, since its only relations modulo the current space are the multiples of $\delta$. So, until it meets a circuit whose dependency lies in one coset, such a scheme gains at most one dimension per circuit and needs at least $a$ circuits to make $\Fp^a$ realizable. If moreover the fresh columns are in general position modulo the current space $U$ (any $a-\dim U$ of them independent modulo $U$), a circuit modulo $U$ has at least $a-\dim U+1$ columns, and a scheme that meets no circuit with dependency in one coset before $\Fp^a$ is realizable uses at least $a(a+3)/2+1=(a+1)(a+2)/2$ columns, the count of Proposition~\ref{prop:triangular-sde}. This is a lower bound for schemes of this form only. Moreover, under the hypotheses of this part, $R_U(S)$ contains a line $\Fp w$, $w\ne0$, if and only if there are $\alpha\in\Fp^{\kappa}$ and a nonconstant map $\beta$ from the set $C=\{\mathrm{cos}(\delta_i):i\in S\}$ to $\Fp$ with $\alpha_\nu+\beta_{c}\in g^{-(\nu+c)}P\cup\{0\}$ for all $\nu\in\mathbb{Z}_\kappa$ and $c\in C$. For $C=\mathbb{Z}_\kappa$ such a pair exists if and only if $p\notin E_2(\kappa)$. So rules that depend on the columns give a circuit a line exactly when a contraction restricted to $C$ does, and at $p\in E_2(\kappa)$ a circuit whose dependency meets all $\kappa$ cosets realizes no line. Sets with a two-dimensional space of relations are not covered: for $\kappa=2$, $p\in\{13,29\}$ and $a=3$, each of four random sets of five columns has $R_0(S)=\Fp^3$ (exhaustive finite computation, AI-checked, pending human review).

\textup{(b)} \emph{Planted hardness below the Chevalley--Warning bound; conditional.} Write $\chi(y)=\mathrm{cos}(y)\in\mathbb{Z}_\kappa$ for $y\ne0$, and let $\theta\in\Fp$ be an unknown key. Suppose an algorithm returns, in time $T$, a nonzero $x\in(P\cup\{0\})^{a+2}$ with $\sum_i x_i\lambda_i=0$ for every $\lambda_1,\dots,\lambda_{a+2}\in\Fp^{a}$ for which such an $x$ exists. Draw $k$ independent uniform points $c_i\in\Fp$ and obtain $\chi(\theta+c_i)$; if some $\theta+c_i$ is $0$, output $-c_i$. Otherwise put $\varepsilon_i=\chi(\theta+c_i)$, pick $u_i\in g^{-\varepsilon_i}P$, and, unless all $c_i$ are equal, take columns in $\Fp^{k-2}$ whose relation space is exactly the span of $u=(u_i)_i$ and $w=(u_ic_i)_i$. The vector $x_i=u_i(\theta+c_i)$ is a solution. Every solution is $\alpha u+\beta w$. If $\beta=0$, all $\varepsilon_i$ are equal. If $\beta\ne0$, then $x_i=\beta u_i(t+c_i)$ with $t=\alpha/\beta$, and for every $i$ either $c_i=-t$ or $\chi(t+c_i)-\chi(\theta+c_i)=-\mathrm{cos}(\beta)$. Since $c\mapsto(t+c)/(\theta+c)$ is a bijection of $\Fp\setminus\{-\theta\}$ onto $\Fp\setminus\{1\}$, for fixed $t\ne \theta$ each $i$ satisfies this with probability at most $1/\kappa+1/(p-1)$, independently of the others. A union bound over $t$ and $\mathrm{cos}(\beta)$ shows that every solution has $\beta\ne0$ and $\alpha/\beta=\theta$, except with probability at most $\kappa^{1-k}+\kappa p\,\kappa^{-k}e^{\kappa k/(p-1)}+(p-1)^{1-k}$. For $k=\lceil\log_\kappa p\rceil+m$ this is at most $\kappa^{1-m}\bigl(1/p+e^{\kappa k/(p-1)}\bigr)+(p-1)^{1-k}$. So, for a constant $m\ge3$ and $p\ge4\kappa k$, $\theta$ is recovered in expected time $T+\poly(k,\log p)$ from an expected $O(k)$ values of $c\mapsto\chi(\theta+c)$ at uniformly random points, each candidate being checked on fresh points. This function is the power-residue pseudorandom function with key $\theta$; for $\kappa=2$ it is the Legendre PRF, the shifted Legendre symbol, whose use as a pseudorandom generator goes back to~\cite{Damgard1988}. No classical key-recovery algorithm polynomial in $\log p$ is known for it or for its power-residue generalization~\cite{BeullensBeyneUdovenkoVitto2020}. Van Dam, Hallgren and Ip~\cite{vanDamHallgrenIp2006} find the shift of the Legendre symbol in quantum polynomial time, but from superposition queries, which this reduction does not supply. Hence, if recovering $\theta$ from $O(\log p)$ values at random points takes classical time superpolynomial in $\log p$ (a cryptographic assumption, not a standard complexity-theoretic one), then no classical algorithm solves every solvable instance with $M=a+2$ columns in time $\poly(a,\log p)$ for primes $p$ of size $\kappa^{a+O(1)}$. For $a\ge2$ these instances have $M=a+2<\kappa a+1$ columns, so the statement concerns solving whenever a solution exists, below the Chevalley--Warning bound. It says nothing about arities at which every instance is solvable, such as $M\ge\kappa a+1$ and, for $p\notin E_2(\kappa)$, the arity of Proposition~\ref{prop:triangular-sde}.

\textup{(c)} For $\kappa=2$ the problem is a homogeneous underdetermined system of $a$ quadratic equations. A polynomial-time algorithm for a class of underdetermined quadratic systems over fields of odd characteristic, with $n\ge m(m+1)/2$ variables for $m$ equations, is due to~\cite{ChengHashimotoMiuraTakagi2014} (threshold as listed in a 2023 survey by Hashimoto; whether it covers every diagonal system is not checked here), so the known thresholds are again quadratic, and recursing through totally isotropic subspaces of one form at a time uses $2^{a+1}-1$ variables. At $p=3$, where $\kappa=2=p-1$, the problem is $\mathbb{F}_3$ subset sum, so a method at arity $o(a^2)$ for $\kappa=2$ that is uniform in $p$ would improve on every known classical polynomial-time algorithm for that problem (Remark~\ref{rem:triangular-scope}).
\end{remark}

\begin{remark}[Steered block-refined triangular solver; conditional; AI-checked, pending human review]
\label{rem:sde-block-steered}
This remark is AI-checked, pending human review (Appendix~\ref{app:provenance}, item~\ref{prov:sde-block-steered}); no human has checked it, and it gives no worst-case improvement. It is adopted, for i.i.d.\ labels only, as the block-steered option of Remark~\ref{rem:combined-schedule} and of rule~(R1) in Appendix~\ref{app:fc-mixed}. Keep the notation of Proposition~\ref{prop:triangular-sde}, with $Q_\nu=g^{\nu}P$ and $Q_\nu^\circ=Q_\nu\cup\{0\}$. Let $S$ be a circuit modulo $U$ whose dependency $\delta$ has entries in at least two cosets, and split $S$ into blocks $S_1,\dots,S_J$, each contained in one set $\{i:\mathrm{cos}(\delta_i)=c_b\}$; put $W_b=\sum_{i\in S_b}\delta_i\lambda_i$. Then the classes of the $W_b$ modulo $U$ span a space $\Omega$ of dimension $J-1$, and their only relation is $\sum_bW_b=e:=\sum_{i\in S}\delta_i\lambda_i\in U$. For $\beta\in\Fp^J$ the weights $x_i=\delta_i(\beta_b+s)$, $i\in S_b$, lie in $P\cup\{0\}$ exactly when $s\in\mathrm{Good}(\beta)=\{s:\beta_b+s\in Q_{-c_b}^\circ\text{ for all }b\}$, and their value is $\sum_b\beta_bW_b+s\,e$; every $x\in(P\cup\{0\})^S$ whose value lies in $\Omega$ modulo $U$ has this form. By the Weil bound~\cite[Theorem~5.41]{LidlNiederreiter1997}, $|\mathrm{Good}(\beta)|\ge\kappa^{-J}\bigl(p-J-(\kappa^J-1)(J-1)\sqrt p\bigr)$ whenever the $\beta_b$ are pairwise distinct. Let $E_{\mathrm{cg}}(J,\kappa)$ be the set of primes $p\equiv1\pmod\kappa$ at which $\mathrm{Good}(\beta)=\emptyset$ for some $J'\le J$, some classes $c_1,\dots,c_{J'}$ not all equal and some pairwise distinct $\beta_1,\dots,\beta_{J'}$, and let $\mathcal{G}_J(p)$ be the least $|\mathrm{Good}(\beta)|$ over such data. Exhaustive computations up to the Weil bound, repeated with independent code (finite computations, AI-checked, pending human review), give $E_{\mathrm{cg}}(3,2)=\emptyset$ and $E_{\mathrm{cg}}(3,3)=\{7,13,19,31,37,43,61,67,73,79,97,103,109,127,139,151,157,163,181,193,211,241\}$. So for $p\notin E_{\mathrm{cg}}(J,\kappa)$ a circuit realises every point of $\Omega$ off the at most $\binom{J}{2}$ hyperplanes $\{\beta_b=\beta_{b'}\}$ with $c_b\ne c_{b'}$, by drawing $s$ uniformly until $J$ coset tests pass, about $p/\mathcal{G}_J(p)$ draws; by part~(a) of Remark~\ref{rem:sde-subquadratic} it never realises a whole plane. Running the triangular recursion with these blocks, and drawing the free scalar $s$ of each round uniformly from $\mathrm{Good}(\beta)$ in the back substitution, moves the next residual along $e$ and steers it off the bad hyperplanes. Suppose that $J\ge\kappa$, that $p\notin E_{\mathrm{cg}}(J,\kappa)$, that the columns of every round, except a last round whose dependency lies in one coset, form a circuit modulo the current space, and that no residual offset difference $\beta_b-\beta_{b'}$ with $c_b\ne c_{b'}$ is identically zero as a linear function of the drawn scalars. The circuit and steering conditions concern the forward pass, including its choice of blocks (the hypothesis $J\ge\kappa$ ensures that every round, whose dependency meets at most $\kappa$ cosets, can be split into at most $J$ blocks), and can be checked in polynomial time once that pass is fixed; when one fails the procedure may abort, and it never returns a wrong solution, since its output is checked. Under these conditions each back substitution returns a nonzero solution with probability at least $1-\Lambda\binom{J}{2}/\mathcal{G}_J(p)$, where $\Lambda\le\lceil a/(J-1)\rceil$ is the number of rounds, and the recursion uses at most
\[
  M_J(a)=1+\mu(a+1)-\tfrac{(J-1)\mu(\mu-1)}{2}+[r_0>0](r_0+1),\qquad a=\mu(J-1)+r_0,\ 0\le r_0<J-1,
\]
columns, which is $(a+1)(a+2)/2$ for $J=2$, $\lfloor(a+2)^2/4\rfloor+[a\text{ odd}]$ for $J=3$, and $a^2/(2(J-1))+O(a)$ in general; the expected time is polynomial in $a$ and $\log p$ times $\kappa^J$. For i.i.d.\ uniform columns both conditions hold with probability at least $1-\bigl(M_J(a)+\Lambda\binom{J}{2}\bigr)/(p-1)$; this is a property of the columns, and repeating the back substitution does not reduce it. Hence, for i.i.d.\ uniform columns in $\Fp^a$ with $a\ge2$ and $J=3$, a randomised polynomial-time procedure with $\lfloor(a+2)^2/4\rfloor+[a\text{ odd}]$ columns fails with probability at most $1/16$ over the columns, and otherwise with probability at most $8^{-k}$ after $k$ back substitutions, for $\kappa=2$ at every prime $p\ge64a^2$ and for $\kappa=3$ at every prime $p\ge\max(2750,1000a^2)$ (both ranges use exhaustive computations of $\mathcal{G}_3(p)$, for $p\le5000$ at $\kappa=2$ and for $p\le11000$ at $\kappa=3$; finite computations, AI-checked, pending human review); with $J=\lceil\log_\kappa a\rceil$ and $p\ge\poly(a)$ it uses $O(a^2/\log a)$ columns. The steering condition is genuinely needed: there are column sets with $M_J(a)\ge\kappa a+1$ columns, hence solvable by the Chevalley--Warning theorem, on which it fails and every back substitution fails (planted instances, exact computations, AI-checked, pending human review: for $\kappa=3$, $J=3$ and $a\in\{8,10\}$, where the blocks are forced to be the coset classes, and for $\kappa=2$, $J=3$ and $a=6$ with a fixed block rule). So for arbitrary columns the best arity proved in this paper remains $(a+1)(a+2)/2$. The labels of Lemma~\ref{lem:s3-law} are i.i.d.\ uniform on the unknown subspace $K_d^\perp$, and the procedure commutes with invertible linear maps of the columns, so run on coordinates in the span of its bank, which is $K_d^\perp$ except with probability at most $(p^{a_d}-1)/((p-1)p^{M})$, it behaves as on i.i.d.\ uniform columns in $\Fp^{a_d}$; with the public arity $M_J(g_d)$ it therefore gives a polynomial-time option that uses the label law, but only at primes large compared with $g_d^2$ and with the number of nodes of the recursion. This is the block-steered option of Remark~\ref{rem:combined-schedule}, used in the combined column of Table~\ref{tab:fcs-budget}. The query counts of Theorem~\ref{thm:main}, of the worst-case rule of Remark~\ref{rem:combined-schedule} and of the other columns of the tables do not use it, and it does not change $E_W(c)$, because its prime condition grows with $g_d$.
\end{remark}

\begin{remark}[Worst case: no quadratic existence barrier; steering by a reserved basis; AI-checked, pending human review]
\label{rem:sde-worst-case}
This remark is AI-checked, pending human review, its computations repeated with independent code (Appendix~\ref{app:provenance}, item~\ref{prov:sde-worst-case}); no human has checked it, and nothing else in the paper uses it. Keep the notation of Remark~\ref{rem:sde-block-steered}, call a sequence of columns a bank, and let $\mathcal{G}^*_J(p)$ be defined as $\mathcal{G}_J(p)$ but allowing all the classes to be equal (and $J'=1$). The Weil bound quoted there gives $\mathcal{G}^*_J(p)\ge\kappa^{-J}\bigl(p-J-(\kappa^J-1)(J-1)\sqrt p\bigr)$. Exact computations up to $p=5000$ and $p=11000$, together with this bound, give $\mathcal{G}^*_3(p)\ge p/10$ for $\kappa=2$ and $p\ge256$, and $\mathcal{G}^*_3(p)\ge p/54$ for $\kappa=3$ and $p\ge2750$; they also show that $\mathcal{G}^*_3(p)\ge1$ holds exactly when $p\ge11$ for $\kappa=2$, and exactly when $p\notin E_{\mathrm{cg}}(3,3)$ for $\kappa=3$ (finite computations, repeated with independent code, AI-checked, pending human review).

\textup{(a)} \emph{Existence is linear in $a$.} Let $N(\kappa,a,p)$ be the largest number of columns in $\Fp^a$ for which~\eqref{eq:sde-target} can have no nonzero solution. Then $2a\le a\,N(\kappa,1,p)\le N(\kappa,a,p)\le\kappa a$. The upper bound is the Chevalley--Warning bound already used in this section. For the lower bound, place unsolvable one-dimensional banks on the coordinates; the two columns $1,-u$ with $u\notin P$ form one. Hence $N(2,a,p)=2a$. Moreover $N(\kappa,1,p)=2$ whenever $\mathcal{G}^*_2(p)\ge1$, for instance whenever $p-2>(\kappa^2-1)\sqrt p$. Indeed, for nonzero $c_1,c_2,c_3$ a solution with third weight $1$ asks for $X\in P\cup\{0\}$ with $X+c_3/c_1\in Q^\circ_{\nu}$, $\nu=-\mathrm{cos}(-c_1/c_2)$, which is a set $\mathrm{Good}(\beta)$ with two distinct offsets. Exhaustive searches give $N(3,2,p)=4$ for $p\in\{13,19,31\}$, $N(3,2,7)=5$ and $N(4,2,17)=5$ (finite computations, repeated with independent code, AI-checked, pending human review). So, for every constant $C_0\ge1$ and every $a>\kappa/C_0$, every bank of $C_0\,a^2$ columns has a nonzero solution, and a quadratic column count can be needed only for finding a solution, not for its existence; for $\kappa=p-1$ this is the zero-sum observation of Theorem~\ref{thm:power-sde}(iii). The gap between the counts of Proposition~\ref{prop:triangular-sde} and Remark~\ref{rem:sde-block-steered} and the existence threshold, at most $\kappa a$, is a search gap. Proving that this gap is necessary for polynomial-time algorithms would be an unconditional lower bound for a total search problem, and is not attempted.

\textup{(b)} \emph{Steering by a reserved basis.} Modify the procedure of Remark~\ref{rem:sde-block-steered} in three ways.
\begin{itemize}
\item Reserve $r$ columns $B$ that form a basis of the span $V$ of the bank, where $r$ is its rank.
\item Build the rounds greedily from the other columns. Scan them, keeping a set $T$ of scanned columns independent modulo the current space $U$, until a column $c$ lies in $U+\mathrm{span}\,T$. The support $S$ of the resulting relation is a circuit modulo $U$; if $|S|\ge2$, split it into $\min(J,|S|)$ blocks, each inside one class of the dependency, and proceed as in Remark~\ref{rem:sde-block-steered}. If $S=\{c\}$, so that $c\in U$, write $c$ in the basis $B$. If $c$ has a nonzero coefficient on some $\sigma\in B\setminus U$, exchange $c$ and $\sigma$ in $B$. Otherwise $c\in\mathrm{span}(B\cap U)$; call this a true-pool event and discard $c$.
\item Give each $\sigma$ in the final $B$ the weight $t_\sigma$, with the $t_\sigma$ independent and uniform in $P$, and start the back substitution from $\tau=-\sum_{\sigma\in B}t_\sigma\sigma$.
\end{itemize}
Every cross-class offset difference $\beta_b-\beta_{b'}$ of every round is then a linear form in $(t_\sigma)$ and the drawn scalars. Its coefficients on the $t_\sigma$ are the values on $B$ of a nonzero functional on $V$, so they are not all zero. Hence neither the circuit condition nor the condition on the offset differences in Remark~\ref{rem:sde-block-steered} is needed.

Define $D_J(0)=0$ and $D_J(q)=\max_{1\le k\le q}\bigl(k+1+D_J(q-\min(J-1,k))\bigr)$. Then $D_3(q)=\lfloor(q+2)^2/4\rfloor$ for $q\ge1$, and $D_J(q)\le q^2/(2(J-1))+q\bigl(1/(J-1)+5/2\bigr)$ for $J\ge3$. Let $J\ge\kappa$ and $\mathcal{G}^*_J(p)\ge1$, and let the bank, of any form, have rank $r$. If the forward pass meets $\pi$ true-pool events and the bank has at least $D_J(r)+r+\pi$ columns, the forward pass reaches $V$. Once it does, each back substitution returns a nonzero solution with probability at least $1-\rho_{\mathrm{pr}}\max\bigl(1/\mathcal{G}^*_J(p),\kappa/(p-1)\bigr)$, and this probability is over the drawn scalars only. Here $\rho_{\mathrm{pr}}\le rJ/2$ is the number of pairs of blocks in different classes, summed over all rounds; for $J=3$, $\rho_{\mathrm{pr}}\le r$ if $\kappa=2$ and $\rho_{\mathrm{pr}}\le3r/2$ if $\kappa=3$.

A true-pool event exhibits a dependent set of at most $r$ columns. So $\pi=0$ for banks in general position, that is, banks in which every $r$ columns span $V$. Hence, for every bank in general position in $\Fp^a$ with at least $\lfloor(a+2)^2/4\rfloor+a$ columns, the procedure with $J=3$ is Las Vegas and fails with probability at most $1/2$ per back substitution, with no randomness over the columns. This holds for $\kappa=2$ at every prime $p\ge\max(256,20a)$, and for $\kappa=3$ at every prime $p\ge\max(2750,162a)$. General position of so many columns already forces $p\ge\lfloor(a+2)^2/4\rfloor+1$ when $r=a\ge2$: the $p+1$ hyperplanes of $V$ through $r-2$ of the columns cover the bank, and each contains at most one further column, so a bank in general position has at most $p+r-1$ columns. Likewise, with $J=\max(\kappa,3,\lceil\log_\kappa a\rceil)$ and $p\ge\max(8\kappa^{2J}J^2,2aJ\kappa^J)$, every bank in general position with $D_J(a)+a$ columns is solved by a Las Vegas procedure in expected time $\kappa^J$ times a polynomial in $a$ and $\log p$. For $a\ge\kappa^\kappa$ this is polynomial time with $O(a^2/\log a)$ columns. On banks of the planted family of Remark~\ref{rem:sde-block-steered} with $a\in\{4,6,8\}$, which are not in general position, padded with uniform columns to $\lfloor(a+2)^2/4\rfloor+a$ columns, the procedure succeeded in all $48$ trials (finite computation, AI-checked, pending human review).

\textup{(c)} \emph{What is open, and the labels.} For arbitrary banks the procedure can stall through true-pool events, and it does so on solvable banks with many proportional columns. It was extended with three repairs, a finishing pass that uses the unused part of $B$, and restarts in random column orders. The repairs are: solving collinear triples by~(a), removing columns outside the span of the others, and recursing on the realised subspace. With them it solved each of $2{,}520$ structured and adversarial banks with $\lfloor(a+2)^2/4\rfloor+a$ columns, for $\kappa\in\{2,3\}$ and $a\le9$ (sampled computation, AI-checked, pending human review), but there is no proof. So for arbitrary columns the best arity proved in this paper remains $(a+1)(a+2)/2$. For the labels of Lemma~\ref{lem:s3-law}, which are i.i.d.\ uniform on $K_d^\perp$, a union bound shows general position except with probability at most $\binom{M}{a_d}/(p-1)$, so this remark applies to them only at primes large compared with $\binom{M}{a_d}$, exponential in $a_d\log a_d$; there Remark~\ref{rem:sde-block-steered} already gives the smaller arity $M_3$ at primes polynomial in the dimension. This remark changes no count of Theorem~\ref{thm:main}, no $E_W(c)$ and no table.
\end{remark}

\begin{remark}[Worst case beyond general position; AI-checked, pending human review]
\label{rem:sde-worst-case-q1}
This remark is AI-checked, pending human review (Appendix~\ref{app:provenance}, item~\ref{prov:sde-worst-case-q1}); no human has checked it, and nothing else in the paper uses it. Keep the notation of Remark~\ref{rem:sde-worst-case}, with $J=3$ and $\kappa\in\{2,3\}$. A round is a set $S$ of columns with a relation vector $\delta$ that is nonzero on $S$, split into blocks each inside one class of $\delta$. Call a family of rounds with disjoint supports a flag if, for each round, the block sums with index $b\ge1$ are independent modulo the space realised by the earlier rounds, and the sum $\sum_{i\in S}\delta_i\lambda_i$ of all its block sums lies in that space. The argument for the success bound in Remark~\ref{rem:sde-worst-case}(b) uses neither the circuit property nor the bound $J$ on the number of blocks, except through $\mathcal{G}^*$. Let $S_0$ be a nonempty set of columns outside the rounds whose span lies in the realised space, give its columns independent weights uniform in $P$, and let $J'$ be the largest number of blocks of a round, with $\mathcal{G}^*_{J'}(p)\ge1$. Suppose that no offset-difference functional of two blocks in different classes vanishes on the span of $S_0$ and of the sums $\sum_{i\in S}\delta_i\lambda_i$ of the later rounds. Then each back substitution started from $S_0$ fails with probability at most $\rho_{\mathrm{pr}}\max\bigl(1/\mathcal{G}^*_{J'}(p),\kappa/(p-1)\bigr)$. Three consequences follow, where $U_L$ is the space realised when the forward pass stops.
(i) If the forward pass stalls with $U_L\ne0$ and $|B\setminus U_L|=\dim V-\dim U_L$, in particular at codimension $1$ with one reserved column outside $U_L$, then the true-pool columns and $B\cap U_L$ span $U_L$, they satisfy this condition, and back substitution on $U_L$ with them as $S_0$ has this failure bound. The condition $U_L\ne0$ holds whenever the bank has at least $D_3(r)+r$ columns.
(ii) Splitting a column $i$ with $\lambda_i\notin U_L$ out of a block of size at least two into its own block yields a flag realising $U_L\oplus\Fp\lambda_i$.
(iii) Suppose a flag realises a hyperplane of $V$, the columns outside its rounds span $V$, and at least three of them lie outside the hyperplane. Then any two of those three form a top round of two singleton blocks, and with the remaining columns outside the rounds as $S_0$ the condition above holds automatically.
With these steps, and with recursion on the columns in the span of $B\cap U_L$, or in the span of the round columns, whenever the counts of the proof guarantee enough columns there, every bank of rank $r$ with at least $\lfloor(r+2)^2/4\rfloor+3r$ columns is reduced in deterministic polynomial time either to a solution, or to a flag together with such a set $S_0$, whose rounds have at most $3+c$ blocks. Here $c\le\lfloor r/2\rfloor$ is the largest number of splittings made in one round, and $c\le2$ unless the splittings at some stall of codimension at least $3$ put three or more into one round. When $c\le2$, the back substitution fails with probability at most $1/2$ for $\kappa=2$ and $p\ge\max(62000,320r)$, and for $\kappa=3$ and $p\ge\max(3.75\cdot10^6,2430r)$. Stalls at which three splittings must go into one round occur for $\kappa\in\{2,3\}$ and $a\in\{7,8\}$. In every one found, the recursion on the span of $B\cap U_L$ applied instead (sampled computation, repeated with independent code, AI-checked, pending human review). No bound on $c$ is proved. So the best arity proved in this paper for arbitrary columns remains $(a+1)(a+2)/2$, and no count of Theorem~\ref{thm:main}, no $E_W(c)$ and no table changes.
\end{remark}

\section{Whole-Grade SDE Schedules}
\label{sec:schedules}
%======================================================================

Each schedule solves~\eqref{eq:sde-target} for whole vector labels in $\Fp^{g_d}$, removing one entire grade per merge.

\subsection{Schedule 1: the Ivanyos--Santha solver}
\label{sec:schedule1}

\begin{theorem}[Ivanyos--Santha {\cite[Theorem~11 and the proof of Theorem~3, Sections~3.1--3.3]{IvanyosSantha2015}}; imported; the reading of the source is AI-checked only (Remark~\ref{rem:is-import}(e)), pending human review]
\label{thm:is-import}
Let $d \ge 1$ be fixed and let $p$ be a prime with $d \mid p-1$. Put $G(d,a) = d^{d(d-1)/2}(a+1)^d$ and $B_d(a) = G(d,a)^{\lceil\log_2(d+1)\rceil}$. There is a deterministic classical algorithm that, given any $\lambda_1,\dots,\lambda_M \in \Fp^a$ with $M \ge B_d(a)$, outputs explicit field elements $s_1,\dots,s_M$, not all zero, with $\sum_{i} s_i^d\lambda_i = 0$, in time polynomial in $M$ and $\log p$, where the polynomial depends on $d$.
\end{theorem}

\begin{remark}
\label{rem:is-import}
\textup{(a)} The algorithm returns the scalars $s_i$ themselves, not only $s_i^d$; Proposition~\ref{prop:descent} needs them, since it applies $\delta_{s_i}$.
\textup{(b)} If $d \nmid p-1$, put $\kappa = \gcd(d,p-1)$ and choose $t \ge 1$ with $t\,d/\kappa \equiv 1 \pmod{(p-1)/\kappa}$. A solution $u$ for the exponent $\kappa$ gives the solution $s_i = u_i^{t}$ for the exponent $d$ \cite[Section~3]{IvanyosSantha2015}, with the same support and in $O(M\log p)$ field operations by square-and-multiply, and $B_\kappa(a) \le B_d(a)$. Indeed $td \equiv \kappa \pmod{p-1}$, so $s_i^d = u_i^{td} = u_i^{\kappa}$ for every $u_i \in \Fp$ (as $t \ge 1$), and $s \ne 0$ when $u \ne 0$; this is Lemma~\ref{lem:app-s2-gcd}(b), checked in Lean as \texttt{ReesSdeGcd} (Section~\ref{sec:formalization}), and~\eqref{eq:s3-bezout} is the same identity, with the Bezout coefficient there (which may be negative, so it is used only on nonzero entries) in the role of $t$. The conversion is therefore proved here, not imported.
\textup{(c)} We use the bound $C(d,a) \le G(d,a)^{\lceil\log_2(d+1)\rceil}$ that the proof of \cite[Theorem~3]{IvanyosSantha2015} establishes. As we read that paper, the bound printed in that theorem, $d^{d^2\log d}(a+1)^{d\log d}$, is not what its proof gives for large $a$, because $d\lceil\log_2(d+1)\rceil > d\log_2 d$; this reading has not been confirmed with the authors or by a human.
\textup{(d)} Determinism is not used in Theorem~\ref{thm:main}: a Las Vegas solver of expected time $T$ would be covered by the solver share of the error budget in Section~\ref{sec:recursion}.
\textup{(e)} This reading of~\cite{IvanyosSantha2015}, including~(a), was checked against arXiv:1503.09016v2 (Appendix~\ref{app:provenance}, item~\ref{prov:sde-is-import}); the reading is AI-checked, pending human review.
\end{remark}

\begin{remark}[Explicit running time of Schedule~1; AI-checked, pending human review; Turing-machine bounds sketched, constants not computed]
\label{rem:is-runtime}
Let $\kappa = \kappa_d$, $a = g_d$, $b = \lceil\log_2 p\rceil$, $\mathcal{G} = G(\kappa,a)$, and, for $\kappa \ge 2$, let $\nu \le b-1$ be the number of prime factors, with multiplicity, of the largest divisor of $p-1$ all of whose prime factors divide $\kappa$ (the case $\kappa_d = 1$ is counted separately below). Let $\pi$ be the set of primes dividing $\kappa$, let $H_\pi$ be the Hall $\pi$-subgroup of the cyclic group $\Fp^\times$, whose order has exactly $\nu$ prime factors counted with multiplicity, and call the component of $y \in \Fp^\times$ in $H_\pi$ its $\pi$-part and the cofactor its $\pi'$-part. The algorithm of Theorem~\ref{thm:is-import} makes at most $\nu+1$ runs, since each restart of \cite[Section~3.1]{IvanyosSantha2015} strictly enlarges a subgroup of $H_\pi$. A deterministic variant, defined for every $\kappa \ge 2$ and described in Remark~\ref{rem:sde-side-tape}, makes at most $r_{\mathrm{run}} = \nu - \min_{q \mid \kappa} v_q\bigl((p-1)/\kappa\bigr)$ runs, with $q$ prime and $v_q$ the $q$-adic valuation; this is $e$ when $\kappa = q^e$ is a prime power, hence a single run, never repeated, when $\kappa$ is prime, and $1 + \max_{q \mid \kappa} v_q(p-1) \le r_{\mathrm{run}} \le \nu$ when $\kappa$ has at least two distinct prime factors (proved; AI-checked, pending human review). Each run makes $\sum_{i<L}\mathcal{G}^i$ calls of \cite[Theorem~11]{IvanyosSantha2015}, with $L = \lceil\log_2 \ell_d\rceil$; counting the field operations of all calls, of the zero sum of $\kappa$-th powers and of the Bezout step of Remark~\ref{rem:is-import}(b) gives at most $40\kappa^3 N_{\mathrm{col}} b(a^2+b^2)$ multiplications and inversions in $\Fp$ and at most twice as many additions, where $N_{\mathrm{col}} = \mathcal{G}^{L} \le B_{\kappa}(a) \le M_d$; at $\kappa_d = 1$ the count is $a^2(a+1)+2a+2b(a+1)$ (Remark~\ref{rem:sde-side-tape}). With schoolbook arithmetic and the binary extended Euclidean algorithm, each operation costs $O(b^2)$ steps of a multitape Turing machine, and all other work is sequential data movement under the tape layout of Remark~\ref{rem:sde-side-tape}. Hence one fixed multitape Turing machine $\mathcal{M}$, with a fixed number of tapes and a fixed alphabet, which reads $d$, $p$, $g_d$, $M_d$ and the columns from its input tape, computes the witness of Schedule~1 within
\[
  T^{\mathrm{IS}}_d = C_1\, d^3\, M_d\, b^3\,(g_d^2 + b^2)
\]
steps, for an absolute constant $C_1$ that is not computed; the machine $\mathcal{M}_d$ of Proposition~\ref{prop:coherent-receiver} is $\mathcal{M}$ with $d$ fixed in its finite control. For fixed $d$ this is\[O\bigl((g_d+1)^{d\lceil\log_2(d+1)\rceil}\, b^3 (g_d^2+b^2)\bigr),\]and $O\bigl(M_d b^2(g_d^2+b)+b^4\bigr)$ when $\nu$ is bounded; the term $b^4$ is the zero sum, which depends only on $p$ and $\kappa_d$. When $\kappa_d = 1$ or $\kappa_d$ is a prime power, in particular for every $d \le 5$, the variant computes a witness of Schedule~1 within
\[
  T^{\mathrm{IS},\ast}_d = C_2\, d^3 (1+\log_2 d)\, M_d\, b^2\,(g_d^2 + b^2) + C_3\, d^4 b^5
\]
steps of one fixed multitape Turing machine, for absolute constants $C_2, C_3$ that are not computed. This is $O\bigl(M_d b^2(g_d^2+b^2) + b^5\bigr)$ for fixed $d$, and the term $C_3 d^4 b^5$ is again the zero sum. Explicit columns attain these run counts (Remark~\ref{rem:sde-side-tape}), so when $\kappa_d$ has at least two distinct prime factors no bound on the number of runs in terms of $\kappa_d$ alone holds for either algorithm, the run count cannot remove the factor $b$, and the bound $T^{\mathrm{IS}}_d$ stands; these are counts of runs, not lower bounds on time (AI-checked, pending human review). The operation counts, the tape layout and the run-count constructions are in Appendix~\ref{app:sde-side-tape}, with provenance in Appendix~\ref{app:provenance}, item~\ref{prov:sde-is-runtime}.
\end{remark}

\paragraph{Arity and determinism (all $p > c$).}
By Theorem~\ref{thm:is-import}, for fixed $d$ the algorithm of Ivanyos and Santha~\cite[Sections~3.1--3.3]{IvanyosSantha2015} returns a nonzero solution of~\eqref{eq:sde-target} for \emph{arbitrary} label columns in $\Fp^{a}$ whenever
\begin{equation}
  \label{eq:is-arity}
  M \ge B_d(a) = \bigl[d^{d(d-1)/2}(a+1)^d\bigr]^{\lceil \log_2(d+1)\rceil},
\end{equation}
in time polynomial in $a$ and $\log p$. The algorithm is deterministic by the proof in Section~3 of~\cite{IvanyosSantha2015}; the statement of Theorem~3 there does not say so. Its Section~3.1 classifies field elements without nonresidues, and its last step uses van de Woestijne's deterministic solver for one equation in $d+1$ variables (over $\Fp$ this is $z_1^d + \cdots + z_{d+1}^d = 0$, equivalently the problem of writing $-1$ as a sum of at most $d$ $d$-th powers; \cite[Section~6]{vandeWoestijne2005}). Only the warm-up Propositions~9 and~10 of~\cite{IvanyosSantha2015} are randomized. Its solver term in the budget of Section~\ref{sec:recursion} is therefore zero. If $d \nmid p-1$, solve with exponent $\kappa_d = \gcd(d,p-1)$ and convert by a Bezout exponent; the displayed $d$-bound remains sufficient. The smaller arity $B_{\kappa_d}(a)$, or $a+1$ when $\kappa_d = 1$, is also sufficient, and it is the arity the reduced solver actually uses; the counts below that use $B_d$ are therefore upper bounds, not exact counts. More precisely, the deterministic run needs only $G(\kappa_d,a)^{\lceil\log_2 \ell_d\rceil}$ columns, with $G$ as in Theorem~\ref{thm:is-import} and $\ell_d \le \kappa_d+1$ the number of terms of the computed zero sum of $\kappa_d$-th powers. A deterministic implementation of the reduced solver (Appendix~\ref{app:provenance}, item~\ref{prov:sde-is-solver}) returns exactly verified witnesses at $\kappa_d \le 3$ at full arity, and at $\kappa_d = 4$ when the computed zero sum has two terms; larger $\kappa_d$ exceed $10^8$ columns already at $a = 1$. It is a classical solver only: no proof in this paper depends on it, although the finite end-to-end simulations and the measured-label circuit backends of Section~\ref{sec:uniform-compilation} call it; it has had no human review, and the general SDE step is not counted as implemented. Every returned witness is checked exactly, and no field enumeration occurs. The bound~\eqref{eq:is-arity} is the one the proof gives (checked against arXiv:1503.09016v2). Section, theorem and corollary numbers for~\cite{IvanyosSantha2015} refer to that arXiv version; the numbering of the journal version cited in the bibliography and of the FAW~2015 extended abstract was not checked. Theorem~11 there finds colliding representations from $G(d,a) = d^{d(d-1)/2}(a+1)^d$ vectors. Section~3.3 iterates it $\lceil\log_2(d+1)\rceil$ times and concludes $C(d,a) \le G(d,a)^{\lceil\log_2(d+1)\rceil}$. Its proof closes with the chain $C(d,a) \le G(d,a)^{\lceil\log_2(d+1)\rceil} \le d^{d^2\log d}(a+1)^{d\log d}$, and its Theorem~3 prints the last bound, as $n > d^{d^2\log d}(m+1)^{d\log d}$ for $n$ variables and $m$ equations (arXiv:1503.09016v2, p.~3). The second inequality fails for every $d \ge 2$ once $a$ is large, because $d\lceil\log_2(d+1)\rceil > d\log_2 d$; for $d = 2$ the exponents of $a+1$ are $4$ and $2$. The first inequality is what the proof gives, and we use it as~\eqref{eq:is-arity}. With $M_d = B_d(g_d)$, the number of lifted sources per terminal sample is $\prod_{d=2}^c B_d(g_d)$. At fixed $c$, the total query count is polynomial in $m$ with exponent at most $1 + \sum_{d=2}^c d\lceil\log_2(d+1)\rceil$; this is $\ReesClassNineQueryExp$ at $c = \ReesClassNineClass$. This is polynomial only for fixed class.
% NEEDS HUMAN VERIFICATION (referee 9, 2026-11-04): the ISSAC 2005 paper has six pages (pp. 348--353); check whether "Section~6" / "Section~6.1" of vandeWoestijne2005 refers to it or to the thesis of the same title.

\subsection{Schedule 2: two-level symmetric contraction}
\label{sec:schedule2}

% E2-UPDATE (1): kappa = 8, 9 sentence at the end of this paragraph ("has not reached that threshold ...
% only [3,20000] remains open ... 1373 listed primes"). E_2(8) is now known exactly; rewrite with final numbers.
\noindent\emph{Terms.} The following objects are defined in full in Appendix~\ref{app:schedule2}; we recall them so that the statements below can be read without it. Fix $\kappa = \kappa_d \ge 2$, put $h = (p-1)/\kappa$, and let $Q_0 = \{x^\kappa : x \in \Fp^\times\}$, the subgroup of index $\kappa$. Fix $\xi \in \Fp^\times$ whose image generates $\Fp^\times/Q_0$, and put $z_m = \xi^m$, $Q_m = z_mQ_0$ and $Q_m^\circ = Q_m \cup \{0\}$ for $m \in \mathbb{Z}/\kappa$. A \emph{coset test} computes the index $m$ with $x \in Q_m$ from $x^h$. The \emph{two-level construction} (Appendix~\ref{app:s2-tensor}) splits $M_d^{\mathrm{sym}} = (a+1)\bigl((\kappa-1)a+1\bigr)$ columns into $(\kappa-1)a+1$ blocks of $a+1$. It takes a linear dependence inside each block, sorts its coefficients by coset, and then takes a dependence among the blocks. The result is a $\kappa\times\kappa$ array of disjoint sets of columns, each with $\kappa$-th-power weights, whose values $T_{ij} \in \Fp^a$ have every coordinate in $W \otimes W$, where $W = \{w \in \Fp^\kappa : \sum_m z_mw_m = 0\}$ (Lemma~\ref{lem:app-s2-tensor}). A \emph{two-level contraction} at $(p,\kappa)$ is a matrix $C \in \Fp^{\kappa\times\kappa}$ with all entries in $Q_0^\circ$, with $C_{00} \ne 0$, and with $\sum_{i,j}C_{ij}X_{ij} = 0$ for every $X \in W\otimes W$. It depends only on $(p,\kappa)$, not on the columns, and in this sense it is universal. Scaling the cells of the array by $\kappa$-th roots of the $C_{ij}$ turns any array into a solution of~\eqref{eq:sde-target} (Lemma~\ref{lem:app-s2-witness}). Every such $C$ has the \emph{decoupled} form $C_{ij} = z_iz_j(\alpha_j+\beta_i)$ with $\alpha,\beta \in \Fp^\kappa$ (Lemma~\ref{lem:app-s2-normal}). The \emph{symmetric} form is $\alpha = \beta$. Its scalars are drawn one at a time by rejection, each against the coset conditions fixed by the earlier ones (Phase~0 of Appendix~\ref{app:s2-symmetric}). Theorem~\ref{thm:power-sde}(ii) uses the symmetric form. $E_2(\kappa)$ is the set of primes $p \equiv 1 \pmod\kappa$ at which no two-level contraction exists. For $\kappa \le 7$ the values of $E_2(\kappa)$ quoted below come from an exhaustive search over $(\beta,\alpha)$ that tests all $\kappa^2$ coset conditions directly, up to an explicit threshold above which a contraction is proved to exist (Corollary~\ref{cor:app-s2-small}; Appendix~\ref{app:provenance}, item~\ref{prov:sde-decoupled}). For $\kappa = 8$ and $9$ the exhaustive search has not reached that threshold, and only partial information is quoted; for $\kappa = 8$, explicit contractions cover every prime in $(20000, \lceil T_A(8)\rceil]$, so only $[3,20000]$ remains open; for $\kappa = 9$ they cover every prime in $(20000, 9^9/2]$ except $1373$ listed primes, and $9^9/2 < T_A(9)$ (see the paragraph after Proposition~\ref{prop:decoupled-sde}). The thresholds $T_A(\kappa)$ and $T_S(\kappa)$ are defined in Appendix~\ref{app:s2-TA} from a closed-form bound $U_\kappa(p)$, decreasing in $p$, on the normalised second moment $\mathbb{E}[n_j^2]/\mathbb{E}[n_j]^2$ of the number $n_j$ of admissible values of $\alpha_j$ when $\beta$ is uniform with $\beta_0 = 0$. For $p > T_A(\kappa)$ a contraction exists, and for $p > T_S(\kappa)$ each round of the sampler of Proposition~\ref{prop:app-s2-TA}(iii) finds one with probability at least $3/8$. Table~\ref{tab:app-s2-thresholds} lists both for small $\kappa$. All of this is AI-checked, pending human review.

\begin{theorem}[Two-level symmetric SDE contraction; part~(ii) conditional, AI-checked, pending human review; proof in Appendix~\ref{app:schedule2}]
\label{thm:power-sde}
Let $\kappa = \kappa_d = \gcd(d,p-1)$ and let the label columns lie in $\Fp^{a}$.
\begin{enumerate}[label=(\roman*),leftmargin=2em]
  \item If $\kappa = 1$, then $a + 1$ columns suffice, since the equation is linear in $s_i^d$ and $s \mapsto s^d$ is a bijection.
  \item If $\kappa \ge 2$ and $p \ge 16\,\kappa^{2\kappa+2}$, the two-level construction with the symmetric form of the contraction (defined above) takes
  \begin{equation}
    \label{eq:power-arity}
    M_d^{\mathrm{sym}} = (a+1)\bigl((\kappa-1)a+1\bigr)
  \end{equation}
  arbitrary columns and an abort parameter $k \ge 1$. It either aborts, with probability at most $2^{-k}$, or returns a nonzero solution, which is checked exactly before output and is never wrong. Its classical cost is at most $(1+\kappa^\kappa)\cdot\poly(a,\kappa,\log p,k)$.
  \item If $\kappa = p-1$, every nonzero $d$-th power equals one, and~\eqref{eq:sde-target} becomes a zero-sum subset problem. Some instances then need many columns. For example, at $p = 5$, $d = 4$, $a = 2$, the columns $e_1,e_1,e_1,e_1,e_2,e_2,e_2,e_2$ have no nonempty zero-sum subset. Nine columns always suffice, since the Davenport constant of $C_5^2$ is $2\cdot 5-1 = 9$ (Olson's theorem for abelian $p$-groups~\cite{Olson1969}). This is an existence statement: Schedule~1 allots $B_4(2) = G(4,2)^3 \approx 3.7\cdot 10^{16}$ columns there, so the small-$p$ regime is a problem of finding cost, not of existence.
\end{enumerate}
\end{theorem}

Part~(ii) is proved in Appendix~\ref{app:schedule2} from the two-level construction and the randomness budget; the proof is AI-checked, pending human review, and the Weil-count step is the one most in need of human review (Appendix~\ref{app:provenance}, item~\ref{prov:sde-power}). The appendix also shows that the proof needs only $p \ge 4\kappa^{2\kappa+2}$ (Remark~\ref{rem:app-s2-four}); the theorem keeps $16\kappa^{2\kappa+2}$. The whole-grade composition and the counts of Table~\ref{tab:full-grade-budget} follow from the arity~\eqref{eq:power-arity}, and a finite computation recomputes every count of that table from the formulas (item~\ref{prov:sde-power}). For every strictly decreasing profile $r = g_2 > g_3 > \cdots > g_c \ge 1$, with $r = \dim[\mathcal{L},\mathcal{L}]$ and $M_d^{\mathrm{sym}}$ taken at $a = g_d$, the bounds $(\kappa_d-1)g_d+1 \le (d-1)(g_d+1)$ and $g_d \le r-d+2$ give $\prod_{d=2}^c M_d^{\mathrm{sym}} \le (c-1)!\prod_{d=2}^c(r-d+3)^2 = (c-1)!\,[(r+1)!/(r-c+2)!]^2 \le r!\,((r+1)!)^2$, since $c-1 \le r$; a finite computation verifies this exactly for $r \le 14$. No human has read the proof, so part~(ii) stays conditional, and Theorem~\ref{thm:main} does not use it. Since $\kappa_d \le c$, the single condition $p \ge 16c^{2c+2}$ covers every grade. A finite computation checks~\eqref{eq:power-arity} with exact witnesses in four cases (item~\ref{prov:sde-power}). On a $61$-bit prime with reduced degree $\kappa = \PsdeDegA$, rank $\PsdeRankA$ uses $\PsdeColsA$ columns and rank $\PsdeRankB$ uses $\PsdeColsB$. On a $128$-bit prime with $\kappa = \PsdeDegC$, rank $\PsdeRankC$ uses $\PsdeColsC$ columns. At $p = \PsdePrimeD$ with $\kappa = \PsdeDegD$, rank $\PsdeRankD$ uses $\PsdeColsD$ columns.

\begin{proposition}[Decoupled two-level contraction; conditional; AI-checked, pending human review]
\label{prop:decoupled-sde}
Let $\kappa = \kappa_d \ge 2$ and let the label columns lie in $\Fp^{a}$. Let $E_2(\kappa)$ be the set of primes $p \equiv 1 \pmod\kappa$ at which no universal two-level contraction of the decoupled form $C_{ij} = z_iz_j(\alpha_j + \beta_i)$ exists, equivalently (Lemma~\ref{lem:app-s2-normal}) no two-level contraction. If $p \notin E_2(\kappa)$, which holds in particular when $\kappa \ge 3$ and $p \ge \kappa^{\kappa+1}$, then the two-level decoupled construction returns a nonzero solution of~\eqref{eq:sde-target} from $M_d^{\mathrm{sym}} = (a+1)\bigl((\kappa-1)a+1\bigr)$ arbitrary columns. The scalars depend only on $(p,\kappa)$ and the coset generator $\xi$. For $p > T_S(\kappa)$, they cost $O(\kappa^{\kappa+1}\log\kappa)$ uniform draws per round, that is $O(\kappa^{\kappa+2}\log\kappa)$ coset tests, times $\poly(\log p, k)$ for abort probability $2^{-k}$; below that, a constant-size search for fixed $\kappa$. Computed values give $T_S(\kappa) \le 10\kappa^{\kappa+1}$ (and $\le 9\kappa^{\kappa+1}$ for $\kappa \ge 4$) for $3 \le \kappa \le 40$; no bound for all $\kappa$ is claimed.
\end{proposition}

% E2-UPDATE (2): "not known exactly", decided ranges 16993 / 16633, 14281 in E_2(8), 20107 in E_2(9),
% ratios 1.23 / 0.34 kappa^{(kappa+1)/2}, 1373 exceptions (3 decided, 1370 undecided). Replace with final counts.
The proposition is AI-checked, pending human review (Appendix~\ref{app:provenance}, item~\ref{prov:sde-decoupled}). The exact searches give $E_2(2) = E_2(3) = \emptyset$ and $E_2(4) = \{5,13,17,29,37,53\}$. They give $E_2(5) = \{11,31,41,61,71,101,131,151,181\}$, $E_2(6)$ equal to a set of $53$ primes, the largest $859$, and $E_2(7)$ equal to a set of $56$ primes, the largest $3221$ (exact search over every prime $p \equiv 1 \pmod 7$ below $5\cdot10^6 > T_A(7)$, with $T_A$ as in Proposition~\ref{prop:app-s2-TA}; finite computation, item~\ref{prov:sde-decoupled}); for $\kappa = 8$ and $9$ the exceptional sets are not known exactly, but exhaustive searches (finite computations, AI-checked, pending human review; Appendix~\ref{app:s2-scope}) decide every prime $p \equiv 1 \pmod 8$ up to $16993$ and every prime $p \equiv 1 \pmod 9$ up to $16633$, and show, by separate runs above these ranges, $14281 \in E_2(8)$ and $20107 \in E_2(9)$. So the largest exceptional prime is at least $1.23\,\kappa^{(\kappa+1)/2}$ for $\kappa = 8$ and at least $0.34\,\kappa^{(\kappa+1)/2}$ for $\kappa = 9$; these are lower bounds only, as the search has not been completed beyond these ranges. Explicit contractions give $E_2(8) \cap (49417, \lceil T_A(8)\rceil] = \emptyset$, with $\lceil T_A(8)\rceil = 117441563$, and $E_2(8) \cap (20000, 49417] = \emptyset$ (Appendix~\ref{app:s2-scope}, with sources; finite computations, AI-checked, pending human review). Since $T_A(8)$ is the proved threshold, $E_2(8) \subseteq [3, 20000]$. For $\kappa = 9$, explicit contractions cover every prime $p \equiv 1 \pmod 9$ in $(20000, 9^9/2]$ except $1373$ listed primes, the largest $424549$; three of them, $20071$, $20089$ and $20107$, lie in $E_2(9)$ by complete runs of the same exhaustive search, and the other $1370$ are not shown to lie in $E_2(9)$; since $9^9/2 < T_A(9)$, this does not determine $E_2(9)$ (Appendix~\ref{app:s2-scope}, with sources; finite computation, AI-checked, pending human review). For $\kappa \ge 10$ only bounds on $E_2(\kappa)$ are known. These sets allow zero entries in the contraction, as the definition does; requiring every entry to be nonzero enlarges $E_2(4)$ and $E_2(5)$ (Appendix~\ref{app:s2-scope}). The proofs are a Weil bound and a second moment over $\beta$ with cyclotomic and Jacobi-sum bounds; the second-moment part was re-derived independently, and neither has had human review. The symmetric form $\alpha = \beta$ used in Theorem~\ref{thm:power-sde}(ii) is strictly weaker. The hypothesis of Theorem~\ref{thm:power-sde}(ii) could be replaced by that of Proposition~\ref{prop:decoupled-sde}; with it, $p \ge c^{c+1}$ would cover every grade with $\kappa_d \ge 3$, and grades with $\kappa_d = 2$ would need no prime condition. We do not adopt the replacement here: Theorem~\ref{thm:power-sde}, Remark~\ref{rem:combined-schedule} and Table~\ref{tab:main-regimes} keep the hypothesis $p \ge 16\,\kappa^{2\kappa+2}$.

\begin{remark}[Size of the two-level threshold; the central-moment threshold argued; AI-checked, pending human review]
\label{rem:two-level-threshold-size}
% E2-UPDATE (3): "below 1.7 kappa^{(kappa+1)/2} for kappa <= 7"; decide whether to extend to kappa = 8 once E_2(8) is final.
It is argued, not proved here, that replacing the second moment over $\beta$ by central moments of order $2m$, with the factorial moments of $|X_j(\beta)|$ controlled by the Weil bound, gives $p \notin E_2(\kappa)$ for every prime $p \equiv 1 \pmod\kappa$ with $p \ge 4e\lceil\ln 4\kappa\rceil^2\kappa^\kappa + 1$ when $\kappa \ge 32$, and the computable form of the same bound improves on $T_A(\kappa)$ for every $\kappa \ge 15$ (Appendix~\ref{app:provenance}, item~\ref{prov:sde-threshold}; status: argued). If that argument is right, the two-level threshold is at most $\kappa^{\kappa+o(1)}$, and combining it with Weil-controlled joint moments of several column counts lowers this to $6\kappa^\kappa + 1$ for $\kappa \ge 729$, an $O(\kappa^\kappa)$ threshold. This method cannot push below $\kappa^\kappa/2$: for $\kappa \ge 3$ and $p \le \kappa^\kappa/2$, an explicit law of the column counts agrees with independent copies on every $\kappa-1$ columns, satisfies the same Weil-controlled moment brackets and gives probability $0$ to every column being nonempty, so no lower bound valid for every law satisfying those brackets is positive (Proposition~\ref{prop:app-s2-barrier}; proved, AI-checked, pending human review). The square gap itself remains open. For the small $\kappa$ of this paper, $T_A$ remains the better explicit bound. The exceptional primes found by exhaustive search lie below $1.7\,\kappa^{(\kappa+1)/2}$ for $\kappa \le 7$, and $\kappa^{(\kappa+1)/2}$ is where the expected number of contractions crosses one. Closing this square gap is open. For composite $\kappa$ it can be bypassed at a cost in arity (Appendix~\ref{app:s2-tower}; AI-checked, pending human review): refining the cosets one prime factor $q_1 \ge \dots \ge q_L$ of $\kappa$ at a time gives a solution of~\eqref{eq:sde-target} from $(a+1)\prod_{t=1}^{L}((q_t-1)a+1)$ arbitrary columns for every prime $p \equiv 1 \pmod\kappa$ with $p > \max(4\kappa^{q_1+1}, 4q_1^2\kappa^5)$ and, if $q_1 \ge 5$, $p > T_A(q_1)$; when $\kappa$ is a power of $2$, $p > ((\kappa-1)^2+1)^2$ suffices. This arity has degree $L+1 \ge 3$ in $a$, so the construction is an alternative at higher arity whose threshold, $\kappa^{q_1+O(1)}$, is polynomial in $\kappa$ when the largest prime factor $q_1$ of $\kappa$ is bounded: it does not lower the Schedule~2 arity, and it changes neither $E_2(\kappa)$ nor $E_W(c)$ nor Table~\ref{tab:main-regimes}. For every $\kappa \ge 4$, prime or not, the gap can also be bypassed at cubic arity (Appendix~\ref{app:s2-chains}; AI-checked, pending human review): a chain of two coset refinements whose first step only asks each entry to lie in one of $k$ cosets gives a solution from $\binom{a+3}{3}$ arbitrary columns for every prime $p \equiv 1 \pmod\kappa$ above an explicit closed-form threshold (Corollary~\ref{cor:app-s2-union-chain}) that is at most $4.05\,\kappa^2(\ln\kappa)^\kappa = \kappa^{\kappa\ln\ln\kappa/\ln\kappa+O(1)}$ for $\kappa \ge 16$; at $\kappa = 5$ and $7$ that closed form, minimised over $k$, is about $2.2\cdot10^5$ and $1.7\cdot10^6$, and evaluating the underlying second-moment criteria numerically instead (Remark~\ref{rem:app-s2-chains-scope}) gives thresholds of about $3781$ and $30823$, against $T_A(\kappa) > (\kappa-1)\kappa^\kappa$. Like the tower, for large $a$ this does not lower the Schedule~2 arity, and it changes neither $E_2(\kappa)$ nor $E_W(c)$ nor Table~\ref{tab:main-regimes}. Barriers to closing the square gap by the standard tools, argued in the cited sources, are collected in Remark~\ref{rem:sde-side-barrier} (Appendix~\ref{app:sde-side-barrier}).
\end{remark}

\begin{proposition}[Triangular two-level solver; AI-checked, pending human review]
\label{prop:triangular-sde}
Let $\kappa=\kappa_d\ge 2$, $p\notin E_2(\kappa)$, and let $(\alpha,\beta)$ be a universal two-level contraction as in Proposition~\ref{prop:decoupled-sde}, normalised by $\alpha_0=1$, $\beta_0=0$ (by translation and by scaling with $\alpha_0^{-1}\in P$; Appendix~\ref{app:s2-contraction}), so that $\alpha_j+\beta_i\in g^{-(i+j)}P\cup\{0\}$ for all $i,j\in\mathbb{Z}_\kappa$, where $P=Q_0=(\Fp^\times)^\kappa$ (the $P_\kappa$ of Schedule~3, not the product $P=\prod_d M_d$ of the budget count) and $g$ is a coset generator (with $g=\xi$, $g^{\nu}P=Q_\nu$ and $\mathrm{cos}=\operatorname{ind}$ of Appendix~\ref{app:s2-gcd}). Then from any
\[
  M_d^{\mathrm{tri}}=\tfrac{(a+1)(a+2)}{2}
\]
columns $\lambda_1,\dots,\lambda_M\in\Fp^{a}$, an algorithm that is given $g$ and $(\alpha,\beta)$ and uses $O(a+1)$ Gaussian eliminations and $O((a+1)^2)$ evaluations of $y\mapsto y^{(p-1)/\kappa}$ returns, deterministically, a nonzero $x \in (P\cup\{0\})^M$ with $\sum_i x_i\lambda_i = 0$. Extracting $\kappa$-th roots of the $x_i$ (Las Vegas, expected time $\poly(\kappa,\log p)$ each) and applying a Bezout exponent gives a nonzero solution of~\eqref{eq:sde-target}. Only the coset $gP$ is used, so for $\kappa = 2$ any quadratic non-residue can serve as $g$; when $g$ is not given, finding one is Las Vegas. The whole procedure is therefore deterministic given $g$, $(\alpha,\beta)$ and the roots, and Las Vegas otherwise; the cost of finding $(\alpha,\beta)$ is that of Proposition~\ref{prop:decoupled-sde}.
\end{proposition}

\begin{proof}
Write $\mathrm{cos}(y)=\nu$ when $y\in g^{\nu}P$. For a column set $S$, a vector $\delta\in\Fp^S\setminus\{0\}$ and a subspace $U\le\Fp^a$ with $e=\sum_{i\in S}\delta_i\lambda_i\in U$, put $v=\sum_{i\in S,\ \delta_i\ne0}\beta_{\mathrm{cos}(\delta_i)}\delta_i\lambda_i$. For $\zeta\in\Fp^\times$ with $\mathrm{cos}(\zeta)=n$, the weights $x_i=\zeta\,\delta_i(\alpha_n+\beta_{\mathrm{cos}(\delta_i)})$ for $\delta_i\ne0$ and $x_i=0$ for $\delta_i=0$ lie in $P\cup\{0\}$, because $\zeta\delta_i\in g^{n+\mathrm{cos}(\delta_i)}P$; and $\sum_i x_i\lambda_i=\zeta\alpha_n e+\zeta v\in \zeta v+U$.

Start with $U_0=0$. In round $k$ take the next $a-\dim U_k+1$ unused columns, choose $\delta\ne 0$ with $e_{k+1}=\sum\delta_i\lambda_i\in U_k$, scale it so that one of its nonzero entries is $1$ (for definiteness the first; Lemmas~\ref{lem:coherent-quadratic-sde} and~\ref{lem:coherent-power-sde} use the last), and form $v_{k+1}$. If $v_{k+1}\in U_k$ (always the case once $U_k=\Fp^a$, when one column is taken with $\delta=1$), stop and let $R$ be the realisation with $\zeta=1$ and $\rho=\sum_iR_i\lambda_i=e_{k+1}+v_{k+1}$ its value; $\rho$ lies in $U_k$, and $R$ has the weight $\alpha_0+\beta_0=1$ on one column. Otherwise set $U_{k+1}=U_k\oplus\Fp v_{k+1}$. Since $v_{k+1}\notin U_k$ in that case, $\dim U_k=k$, so round $k$ takes $a-k+1$ columns, and round $a$, if reached, takes one column and stops. At most $\sum_{k=0}^{a}(a-k+1)=(a+1)(a+2)/2$ columns are used.

Put $\tau=-\rho$. For $l=k,\dots,1$ let $\zeta_l$ be the $v_l$-coordinate of $\tau\in U_{l-1}\oplus\Fp v_l$; if $\zeta_l\neq0$, add the weights of round $l-1$ (the round that formed $e_l$ and $v_l$) for $\zeta=\zeta_l$ and replace $\tau$ by $\tau-\zeta_l(\alpha_{\mathrm{cos}(\zeta_l)}e_l+v_l)$, which lies in $U_{l-1}$. Each update subtracts the value of the weights just added, so at the end $\tau$ equals minus the value of all the weights, and $\tau\in U_0=0$. The supports are disjoint, so the weights sum to a nonzero solution of $\sum x_i\lambda_i=0$ with $x_i\in P\cup\{0\}$. Taking $\kappa$-th roots and a Bezout exponent gives a solution of~\eqref{eq:sde-target} by Lemma~\ref{lem:app-s2-gcd}(b).

Cost. Each round uses one elimination for $\delta$ and one for the test $v_{k+1}\in U_k$, and the back substitution one coordinate solve per round, so there are at most $3(a+1)$ eliminations, each of size at most $a\times(a+1)$. Coset tests are made for the at most $M$ nonzero $\delta_i$ and the at most $a$ values $\zeta_l$. With $\delta$ chosen by a fixed rule, as above, the procedure is deterministic given $g$ and $(\alpha,\beta)$. Any $g'\in gP$ gives the same run with the same $(\alpha,\beta)$; a generator of another coset needs the transported contraction of Appendix~\ref{app:s2-contraction}. A finite computation (AI-checked, pending human review; Appendix~\ref{app:provenance}, item~\ref{prov:sde-triangular}) ran the procedure on $178{,}080$ random and adversarial column sets for $\kappa\in\{2,3,4\}$, $p<10^4$ and $a\le 6$, and exhaustively on $603{,}397$ small tuples, with no failure.
\end{proof}

\begin{remark}[Scope and priority of the triangular solver]
\label{rem:triangular-scope}
The proposition is AI-checked, pending human review (Appendix~\ref{app:provenance}, item~\ref{prov:sde-triangular}); no human has checked it, and it does not change Theorem~\ref{thm:power-sde}, Remark~\ref{rem:combined-schedule} or Table~\ref{tab:main-regimes}. Since $p \ge 16\kappa^{2\kappa+2}$ implies $p \notin E_2(\kappa)$ by the proof of Theorem~\ref{thm:power-sde}(ii) (conditional; see the paragraph after that theorem), it would lower the Schedule~2 arity of the combined rule to $(g_d+1)(g_d+2)/2$ under the same guard; the query counts of Theorem~\ref{thm:main}, Remark~\ref{rem:combined-schedule} and Tables~\ref{tab:main-regimes}, \ref{tab:fcs-budget} and~\ref{tab:query-crossover} do not use that (the coherent-label circuits of Section~\ref{sec:uniform-compilation} and Lemmas~\ref{lem:coherent-quadratic-sde} and~\ref{lem:coherent-power-sde} run at that arity, and Remark~\ref{rem:fc-tradeoff}(a) uses its rounds). Adopting it there would not change $E_W(c)$ for $2\le c\le 9$: both arities are quadratic in $g_d$, and the maximising primes are governed by Schedule~1 grades. It would lower only constants; for example, at $p=2^{61}-1$ it would lower the worst-case combined count given in the caption of Table~\ref{tab:fcs-budget} from $2^{555.4}$ to $2^{542.0}$ (finite computation, item~\ref{prov:sde-triangular}), while the block-steered option of Remark~\ref{rem:combined-schedule}, which uses the label law, gives $2^{533.4}$ there. It is therefore not adopted here. For $\kappa=2$ the proposition is the dimension-reduction bound of~\cite[Claim~1]{IvanyosSanselmeSantha2012} for diagonal quadratic equations; at $p=3$ it is the $\mathbb{F}_3$ subset-sum bound recorded in~\cite[Theorem~2.2]{KothariODonnellWu2026}. Using sparse dependencies~\cite[Lemma~3.6]{KothariODonnellWu2026} lowers the count to $(1-1/p)a^2/2+O(a\log a)$. It is argued, not proved here (item~\ref{prov:sde-triangular}; status: argued), that the argument applied to Schedule~2 lowers its second level to $a+1$ objects, giving $(a+1)^2$ columns in place of $(a+1)((\kappa-1)a+1)$, and that schemes built from a universal contraction, in the sense of that source, cannot go below $(a+1)^2$ and $(a+1)(a+2)/2$, because every such contraction has at least two levels with a nontrivial relation. No polynomial-time method at arity $o(a^2)$ is known for arbitrary columns and any fixed $\kappa\ge2$ (Appendix~\ref{app:sde-side-arity}: Remark~\ref{rem:sde-block-steered} gives one for i.i.d.\ uniform columns, and Remark~\ref{rem:sde-worst-case} one for every bank in general position, both at primes $p \ge \poly(a)$). Remark~\ref{rem:sde-subquadratic} (AI-checked, pending human review) shows that, for columns in general position, schemes that enlarge the realizable space one circuit at a time and meet no circuit whose dependency lies in one coset need $(a+1)(a+2)/2$ columns, and that, under a cryptographic assumption, no classical polynomial-time algorithm solves every solvable instance at arity $a+2$ for primes of size $\kappa^{a+O(1)}$; neither statement excludes a method at an arity at which every instance is solvable. At $\kappa=p-1$ the problem is the problem $\mathrm{BIS}_p$ of~\cite[Definition~6.1]{GoosKamathSotirakiZampetakis2020}, for which every known polynomial-time algorithm needs a number of columns at least quadratic in $a$~\cite{KothariODonnellWu2026,KothariMetgerODonnellShuttyWu2026}.
% NEEDS HUMAN VERIFICATION: Theorem 2.2 and Lemma 3.6 of KothariODonnellWu2026 -- numbering of arXiv:2510.07515 or of the STOC 2026 pages 101--105? Also "arXiv version, Section 5.2" in the introduction (referee 12, 2026-11-04).
\end{remark}

\subsection{The combined rule}
\label{sec:combined-rule}

\begin{remark}[Combined per-grade schedule; conditional]
\label{rem:combined-schedule}
The schedule can be chosen grade by grade from public data only, namely $p$, $d$ and $g_d$, never from $H$, $K_d$ or $a_d$. For each $d = 2,\dots,c$: if $\kappa_d = 1$, use the linear schedule of Theorem~\ref{thm:power-sde}(i) with $M_d = g_d+1$; if $\kappa_d \ge 2$ and $p \ge 16\,\kappa_d^{2\kappa_d+2}$, use Schedule~2 with $M_d = (g_d+1)\bigl((\kappa_d-1)g_d+1\bigr)$; otherwise use Schedule~1 at degree $\kappa_d$ with $M_d = B_{\kappa_d}(g_d) \le B_d(g_d)$, converted by a Bezout exponent. Because Schedule~1 is the fallback, this rule is valid for every prime $p > c$ and every $H$, and it uses no label law. The recursion of Section~\ref{sec:recursion} with these arities returns $\hh$ with probability at least $1-\eps$, conditional on the steps behind Theorem~\ref{thm:main} and on part~(ii) above (AI-checked, pending human review; Appendix~\ref{app:provenance}, item~\ref{prov:sde-combined}); no human has checked it. Most of its saving at small $p$ comes from running Schedule~1 at the reduced degree $\kappa_d$, not from Schedule~2. The worst case over primes barely moves: at fixed $c$, the exponent of $m$ in the query count, maximized over all $p > c$, is unchanged for $c \le 7$, and it falls by two for $c = 8$ and $c = 9$, to $143$ at $c = 9$ against the bound $\ReesClassNineQueryExp$ of Schedule~1. The rule itself uses no label law and does not mix in Schedule~3. Proposition~\ref{prop:mixed-schedule} in Appendix~\ref{app:s3-mixed} mixes the uniform-kernel nodes of Theorem~\ref{thm:free-coordinate}(iv) into it, with a separate $u_d$ at each grade, the node-count-weighted bank bound $\sum_d n_d\beta_d \le 1/64$ and the tree restarts of Theorem~\ref{thm:free-coordinate}(ii) and~(iv); it needs the label law at those grades, and Theorem~\ref{thm:power-sde}(ii) at any Schedule~2 grade it keeps. Theorem~\ref{thm:free-coordinate}(iii) below gives a variant that does mix Schedule~3 into this rule at every grade with $2 \le \kappa_d < p-1$, with node-count-weighted bank bounds and one union bound; it needs the label law at those grades.

\emph{Block-steered option (uses the label law; AI-checked, pending human review).} At a grade with $2 \le \kappa_d < p-1$ the rule may instead use the block-steered node of Remark~\ref{rem:sde-block-steered}, with $J_d = \max(3,\kappa_d)$ blocks and $M_d = M_{J_d}(g_d)$ columns, when this arity is smaller and three public conditions on $p$, $d$, $g_d$ and the earlier arities hold. Put $\Lambda_d = \lceil g_d/(J_d-1)\rceil$ and $G^-_d = \min_{2\le j\le J_d}\kappa_d^{-j}\bigl(p-j-(\kappa_d^{\,j}-1)(j-1)\sqrt p\bigr)$, the Weil bound of that remark. (G1)~$p \notin E_{\mathrm{cg}}(J_d,\kappa_d)$; for example $G^-_d > 0$, which makes every $\mathrm{Good}(\beta)$ nonempty. (G2)~$\Lambda_d\binom{J_d}{2} \le G^-_d/2$, so that each back substitution fails with probability at most $1/2$. (G3)~$\bigl(M_d + \Lambda_d\binom{J_d}{2}\bigr)/(p-1) + p^{g_d-M_d}/(p-1) \le \tau_d$, with $\tau_d$, $S$ and $n_d$ as in Theorem~\ref{thm:free-coordinate}(iii) below, so that the bank and span failures of that remark are at most $\tau_d$. The node works in coordinates of the span of its bank, makes at most $B$ back substitutions, with $B$ as in that theorem, and caps its draw loops and root extractions so that each kind of capped step fails with probability at most $2^{-B}$. On the i.i.d.\ supply of Lemma~\ref{lem:s3-law} it therefore aborts with probability at most $\tau_d + 3\cdot2^{-B}$, never returns a wrong witness, converts by a Bezout exponent as in (F4)--(F6) of Appendix~\ref{app:fc-setting}, and runs in time polynomial in $g_d$, $B$ and $\log p$ times $p/G^-_d$. A run with this option is a run of rule~(R1) of Appendix~\ref{app:fc-mixed}, with the error split and the guarantee of Theorem~\ref{thm:free-coordinate}(iii). It needs the label law at the grades that use the node, and Theorem~\ref{thm:power-sde}(ii) only at the Schedule~2 grades it keeps. The combined column of Table~\ref{tab:fcs-budget} uses this option. Condition~(G3) asks $p$ to be large compared with $g_d^2$ and with $n_d$, so at a fixed prime the option is unavailable for large $g_d$, and it would not change $E_W(c)$; in any case Theorems~\ref{thm:main} and~\ref{thm:bb-envelope} and $E_W(c)$ use the rule above without it. The recomputation is recorded in Appendix~\ref{app:provenance}, item~\ref{prov:sde-combined}.
\end{remark}

